\subsection{Composition, Restriction, and Inverse Functions}

\begin{definition} \label{def:composition-restriction-extension-identity} \hlt{Composition, Restriction, Extension, and Identity}\\
Let $X,Y$ be sets and let $f:X\to Y$.
\begin{enumerate}[label=(\roman*)]
\setlength{\itemsep}{0pt}
\item \hlt{Restriction}: If $A\subseteq X$, then $f|_A:A\to Y$ satisfies $(\forall x\in A)\,(f|_A)(x)=f(x)$.
\item \hlt{Extension}: If $U$ is a set with $X\subseteq U$, an extension of $f$ is a map $\widetilde f:U\to Y$ satisfying $\widetilde f|_X=f$.
\item \hlt{Composition}: If $C,Z$ are sets, $g:C\to Z$, and $f(X)\subseteq C$, then $g\circ f:X\to Z$ satisfies $(\forall x\in X)\,(g\circ f)(x)=g(f(x))$.
\item \hlt{Identity}: The map $\operatorname{id}_X:X\to X$ satisfies $(\forall x\in X)\,\operatorname{id}_X(x)=x$.
\end{enumerate}
\end{definition}

\begin{theorem} \label{thm:bijective-iff-invertible-function} \hlt{Bijective Means Invertible}\\
Let $X,Y$ be sets and let $f:X\to Y$.
\begin{enumerate}[label=(\roman*)]
\setlength{\itemsep}{0pt}
\item The map $f$ is bijective iff $(\exists g:Y\to X)\,[g\circ f=\operatorname{id}_X\land f\circ g=\operatorname{id}_Y]$.
\item If $g,h:Y\to X$ both satisfy the two identities in (i), then $g=h$. When it exists, this map is denoted $f^{-1}:Y\to X$.
\end{enumerate}
\begin{proof}
\begin{enumerate}[label=(\roman*)]
\setlength{\itemsep}{0pt}
\item[$(i)$] ($\Rightarrow$) Suppose $f$ is bijective. For $y\in Y$, choose $x_y\in X$ with $f(x_y)=y$ by surjectivity; it is unique by injectivity.\\
Define $g:Y\to X$ by $g(y)\defeq x_y$. For $x\in X$, $g(f(x))=x$, and for $y\in Y$, $f(g(y))=y$.\\
Thus $g\circ f=\operatorname{id}_X$ and $f\circ g=\operatorname{id}_Y$ by function extensionality.
\item[$(i)$] ($\Leftarrow$) Suppose $g:Y\to X$ satisfies the two identities. If $x_1,x_2\in X$ and $f(x_1)=f(x_2)$, then $x_1=g(f(x_1))=g(f(x_2))=x_2$.\\
Hence $f$ is injective. For $y\in Y$, $y=f(g(y))\in\operatorname{range}(f)$, so $f$ is surjective and hence bijective.
\item[$(ii)$] Suppose $g,h:Y\to X$ satisfy the two identities. For $y\in Y$, $g(y)=g(f(h(y)))=h(y)$. Thus $g=h$ by function extensionality.
\end{enumerate}
\hlt{Structural Commentary}: Construction proves existence; cancellation recovers bijectivity and proves uniqueness.\\
\hlt{Key Insight}: A bijection pairs each output with exactly one input.
\end{proof}
\end{theorem}

\begin{proposition} \label{prop:composition-restriction-calculus} \hlt{Composition and Restriction Calculus}\\
Let $X,Y,Z,W$ be sets, and let $f:X\to Y$, $g:Y\to Z$, and $h:Z\to W$.
\begin{enumerate}[label=(\roman*)]
\setlength{\itemsep}{0pt}
\item $h\circ(g\circ f)=(h\circ g)\circ f$.
\item $\operatorname{id}_Y\circ f=f$ and $f\circ\operatorname{id}_X=f$.
\item $(\forall A\subseteq X)\,(g\circ f)(A)=g(f(A))$.
\item $(\forall C\subseteq Z)\,(g\circ f)^{-1}(C)=f^{-1}(g^{-1}(C))$.
\item $(\forall A\subseteq X)\,(g\circ f)|_A=g\circ(f|_A)$.
\item If $f$ and $g$ are bijective, then $g\circ f$ is bijective.
\item If $f$ and $g$ are bijective, then $(g\circ f)^{-1}=f^{-1}\circ g^{-1}$.
\end{enumerate}
\begin{proof}
\begin{enumerate}[label=(\roman*)]
\setlength{\itemsep}{0pt}
\item For $x\in X$, both sides evaluate to $h(g(f(x)))$. Hence the two maps are equal by function extensionality.
\item For $x\in X$, $(\operatorname{id}_Y\circ f)(x)=f(x)=(f\circ\operatorname{id}_X)(x)$. Hence both identities hold by function extensionality.
\item Fix $A\subseteq X$ and $z\in Z$. Then $z\in(g\circ f)(A)\Longleftrightarrow(\exists x\in A)\,z=g(f(x))\Longleftrightarrow(\exists y\in f(A))\,z=g(y)\Longleftrightarrow z\in g(f(A))$.
\item Fix $C\subseteq Z$ and $x\in X$. Then $x\in(g\circ f)^{-1}(C)\Longleftrightarrow g(f(x))\in C\Longleftrightarrow f(x)\in g^{-1}(C)\Longleftrightarrow x\in f^{-1}(g^{-1}(C))$.
\item Fix $A\subseteq X$ and $x\in A$. Then $((g\circ f)|_A)(x)=g(f(x))=(g\circ(f|_A))(x)$, so the restrictions are equal by function extensionality.
\item Suppose $f$ and $g$ are bijective. Fix $x_1,x_2\in X$ and suppose $(g\circ f)(x_1)=(g\circ f)(x_2)$. Then $f(x_1)=f(x_2)$ by the injectivity of $g$. Hence $x_1=x_2$ by the injectivity of $f$, so $g\circ f$ is injective.\\
For $z\in Z$, choose $y\in Y$ with $g(y)=z$ by the surjectivity of $g$, then choose $x\in X$ with $f(x)=y$ by the surjectivity of $f$.\\
Thus $(g\circ f)(x)=z$, so $g\circ f$ is surjective and therefore bijective.
\item Suppose $f$ and $g$ are bijective. The inverse functions $f^{-1}:Y\to X$ and $g^{-1}:Z\to Y$ exist by Theorem \ref{thm:bijective-iff-invertible-function}(i).\\
For $x\in X$, $(f^{-1}\circ g^{-1})((g\circ f)(x))=f^{-1}(g^{-1}(g(f(x))))=f^{-1}(f(x))=x$.\\
For $z\in Z$, $(g\circ f)((f^{-1}\circ g^{-1})(z))=g(f(f^{-1}(g^{-1}(z))))=g(g^{-1}(z))=z$. Thus $(g\circ f)^{-1}=f^{-1}\circ g^{-1}$ by Theorem \ref{thm:bijective-iff-invertible-function}(ii).
\end{enumerate}
\hlt{Structural Commentary}: Pointwise evaluation proves the calculus; injectivity and surjectivity pass through the factors.\\
\hlt{Key Insight}: Inverse maps cancel compositions in reverse order.
\end{proof}
\end{proposition}
